\subsection{Suma en Rango con Actualización en Rango (Dos Fenwick Trees)}
\label{alg:3-5}

\graybold{Preguntas guía:}

\begin{itemize}
\item ¿Debo sumar un valor a TODOS los elementos de un rango, e intercalar consultas de SUMA sobre rangos?
\item ¿La operación de consulta es la suma (invertible), y no mínimo o máximo, que no se pueden "deshacer"?
\item ¿Con $n, q \sim 10^5$ un segment tree con lazy propagation recursivo sería demasiado lento en Python?
\item ¿Puedo expresar la suma de prefijos del arreglo como combinación lineal de dos sumas de prefijos más simples?
\end{itemize}

\graybold{Utilidad:} Resuelve la combinación "sumar $v$ en $[l, r]$" y "consultar la suma de $[l, r]$" en tiempo logarítmico usando solo dos árboles de Fenwick (BIT), sin lazy propagation. Es la alternativa estándar al segment tree con lazy cuando tanto la actualización como la consulta son de suma: el código es más corto, no usa recursión y en Python resulta varias veces más rápido. Para mínimo o máximo en rango con actualización en rango no aplica, porque depende de que la suma sea invertible; ahí sí hace falta lazy propagation.

\graybold{Complejidad:} \bigO{\log n} por actualización y por consulta (dos recorridos de BIT en cada una), \bigO{n} de memoria.

\graybold{Fundamento teórico:} Un BIT guarda en la posición $i$ la suma de un bloque de longitud $i \mathbin{\&} (-i)$ (el bit menos significativo de $i$) que termina en $i$: sumar un prefijo recorre $i \mathrel{-}= i \mathbin{\&} (-i)$ y actualizar un punto recorre $i \mathrel{+}= i \mathbin{\&} (-i)$, ambos en \bigO{\log n} pasos. Por sí solo, un BIT hace actualización PUNTUAL y consulta de PREFIJO. El truco para rangos parte del arreglo de diferencias $d$: sumar $v$ a $[l, r]$ equivale a $d_l \mathrel{+}= v$ y $d_{r+1} \mathrel{-}= v$, y el valor de cada elemento es $a_i = \sum_{j \le i} d_j$. La suma de prefijos que se quiere consultar es entonces una doble suma que se puede reordenar:
\begin{align*}
S(x) &= \sum_{i=1}^{x} a_i = \sum_{i=1}^{x}\sum_{j=1}^{i} d_j = \sum_{j=1}^{x} d_j\,(x - j + 1)\\
     &= x\sum_{j=1}^{x} d_j \;-\; \sum_{j=1}^{x} d_j\,(j-1).
\end{align*}
Ambas sumas son sumas de prefijo de arreglos que solo cambian en DOS puntos por actualización: el primer BIT guarda $d_j$ y el segundo guarda $d_j\,(j-1)$. Así, una actualización de rango son dos actualizaciones puntuales en cada BIT (en $l$ con $v$ y $v(l-1)$, en $r+1$ con $-v$ y $-v\,r$), y la suma de $[l, r]$ es $S(r) - S(l-1)$, con cada $S$ calculado recorriendo ambos BIT a la vez. El factor $x$ se aplica al final, fuera del recorrido, porque multiplica a toda la suma del primer BIT.~\cite{fenwick1994}

\graybold{Ejercicio aplicado}

(SPOJ HORRIBLE — Horrible Queries) Dados $T$ casos, cada uno con un arreglo de $N$ ceros y $C$ comandos que pueden ser sumar $v$ a todas las posiciones de $[p, q]$ o consultar la suma de $[p, q]$, responder todas las consultas.

\graybold{I/O:} $N, C \le 10^5$ por caso, con comandos de 3 o 4 números según el tipo $\Rightarrow$ lectura total + tokenización con índice \texttt{idx} (patrón 2), leyendo primero el tipo del comando para saber cuántos tokens consumir. El enunciado no garantiza $p \le q$, así que se ordenan en cada comando. Las sumas pueden llegar a \aprox{}\ten{17}, sin problema para los enteros de Python. Verificado contra el caso de ejemplo y contrastado con una fuerza bruta (sumar y consultar sobre la lista directamente) en 250 tandas con varios casos cada una e índices con $p > q$, sin discrepancias. Con $N = C = 10^5$ tarda \aprox{}0.2s a \aprox{}0.46s por caso según la mezcla de comandos, mientras que un segment tree con lazy propagation recursivo, verificado igual de correcto, tarda \aprox{}0.75s a \aprox{}2.86s: unas 6 veces más lento. Como $T$ no está acotado, con varios casos máximos el tiempo crece linealmente (\aprox{}2.4s con $T = 5$).

\graybold{Código solución}

\begin{lstlisting}[language=Python]
import sys

def solve():
    data = sys.stdin.buffer.read().split()
    idx = 0
    t = int(data[idx]); idx += 1
    out = []
    for _ in range(t):
        n = int(data[idx]); c = int(data[idx+1]); idx += 2
        b1 = [0]*(n + 1)         # BIT de d[j]
        b2 = [0]*(n + 1)         # BIT de d[j]*(j-1)
        for _ in range(c):
            if data[idx] == b"0":
                p = int(data[idx+1]); q = int(data[idx+2]); v = int(data[idx+3]); idx += 4
                if p > q: p, q = q, p            # el enunciado no garantiza p <= q
                # d[p] += v  y  d[q+1] -= v, reflejado en ambos BIT
                i = p; w = v*(p - 1)
                while i <= n:
                    b1[i] += v; b2[i] += w
                    i += i & -i
                i = q + 1; w = v*q
                while i <= n:
                    b1[i] -= v; b2[i] -= w
                    i += i & -i
            else:
                p = int(data[idx+1]); q = int(data[idx+2]); idx += 3
                if p > q: p, q = q, p
                # S(x) = x * sum(d[1..x]) - sum(d[j]*(j-1))  ->  S(q) - S(p-1)
                s1 = s2 = 0; i = q
                while i > 0:
                    s1 += b1[i]; s2 += b2[i]
                    i -= i & -i
                res = s1*q - s2
                s1 = s2 = 0; i = p - 1
                while i > 0:
                    s1 += b1[i]; s2 += b2[i]
                    i -= i & -i
                res -= s1*(p - 1) - s2
                out.append(res)
    sys.stdout.write("\n".join(map(str, out)) + "\n")

solve()
\end{lstlisting}
